% realtime-webrtc-callkit.tex — real-time beyond WebSockets: WebRTC (signalling vs media, SDP offer/answer,
% ICE, STUN/TURN, NAT traversal, DTLS-SRTP, topologies), CallKit (CXProvider, CXCallController, reporting
% incoming calls), PushKit VoIP pushes + the iOS 13 must-report rule, audio session during calls,
% background behaviour, SignalR / Socket.IO in one line.
% Picture: call setup sequence — caller -> signalling -> VoIP push -> CallKit UI -> answer -> ICE -> media.
% Sources (checked 2026-09-25): developer.apple.com — PKPushRegistryDelegate
% pushRegistry(_:didReceiveIncomingPushWith:for:completion:) (iOS 13 rule, quoted), "Responding to VoIP
% Notifications from PushKit" (apns-expiration 0, fulfil answer after connecting, remoteEnded / failed,
% no PushKit without CallKit -> use an NSE), "Sending notification requests to APNs" (push type voip,
% topic = bundle ID + .voip), CXProviderDelegate provider(_:didActivate:). WebRTC: IETF RFC 8829/9429
% (JSEP), 3264 (offer/answer), 8445 (ICE), 8838 (trickle), 8489 (STUN), 8656 (TURN), 5764 (DTLS-SRTP),
% 7874 / 7742 (Opus; VP8 + H.264). RTCAudioSession useManualAudio / isAudioEnabled: libwebrtc ObjC SDK.
% NOT repeated: WebSocket mechanics (networking/websockets-realtime.tex), APNs + NSE basics
% (push-notifications.tex), categories/modes table (avfoundation-camera-audio.tex), background modes
% list (background-execution.tex).
% @labels: area=ios-platform kind=api platform=ios topic=networking,protocols,platform-apis
% @tags: webrtc, sdp, offer-answer, ice, stun, turn, nat-traversal, dtls-srtp, sfu, callkit, pushkit, voip-push
\documentclass[8pt]{extarticle}
\usepackage{printup-sheet}
\usepackage{array}

\lstdefinelanguage{SwiftSheet}{
  morekeywords={protocol,class,final,struct,enum,func,var,let,weak,init,for,in,
    if,else,return,guard,self,nil,try,await,async,throws,private,some,as,
    true,false,Task,Void,escaping},
  sensitive=true, morecomment=[l]{//}, morestring=[b]",
  literate={->}{{\hbox{-}\hbox{>}}}2 {??}{{\hbox{?}\hbox{?}}}2}
\lstset{basicstyle=\ttfamily\scriptsize, aboveskip=2pt, belowskip=2pt}
\newcommand\ct[1]{\texttt{#1}}
\newcommand\hd[1]{\par\vspace{3pt}\noindent{\bfseries\color{sheetBlue}#1}\par\vspace{1pt}}

\begin{document}

\sheettitle{Real-time calls on iOS — WebRTC, CallKit, PushKit}{ios-platform · folio}

\oneliner{A call is two planes: \textbf{signalling} (your server, any transport — usually a WebSocket —
carrying \textbf{SDP offer/answer} and \textbf{ICE candidates}) and \textbf{media} (WebRTC:
\textbf{ICE} finds a path through NATs with \textbf{STUN}, falls back to a \textbf{TURN} relay, then
\textbf{DTLS-SRTP} encrypts audio/video peer-to-peer). On iOS the app is usually not running when the
call arrives: a \textbf{PushKit VoIP push} wakes it, it \emph{must} report the call to \textbf{CallKit}
at once, and CallKit owns the system call UI and the audio session.}

\vspace{3pt}
\noindent\begin{tikzpicture}[sheet, yscale=1.12,
    lane/.style={box, font=\scriptsize\bfseries, minimum width=21mm, minimum height=5mm, inner sep=1.5pt},
    life/.style={draw=sheetGrey, dashed},
    m/.style={font=\tiny, inner sep=1pt, fill=white, fill opacity=0.9, text opacity=1},
    msg/.style={->, thick, draw=sheetBlue},
    ice/.style={<->, thick, dashed, draw=sheetGrey},
    med/.style={<->, very thick, draw=sheetGreen!70!black},
    act/.style={box, font=\tiny, inner sep=1.2pt, draw=sheetOrange, fill=sheetOrange!10, align=center},
    n/.style={circle, fill=sheetOrange, text=white, font=\tiny\bfseries, inner sep=0.4pt, minimum size=3mm}]
  \foreach \k/\x/\t in {ca/0.9/Caller app, sg/4.3/Signalling server, ap/7.2/APNs (VoIP),
                        ck/10.1/Callee: PushKit + CallKit, cb/13.0/Callee app (WebRTC), st/15.75/{STUN · TURN}} {
    \node[lane] (\k) at (\x,0) {\t};
    \draw[life] (\k.south) -- (\x,-5.05);
  }
  % 1 offer
  \node[act, anchor=west] at (-0.2,-0.6) {\ct{CXStartCallAction} $\to$ \ct{createOffer}};
  \node[n] at (1.1,-1.0) {1};
  \draw[msg] (0.9,-1.0) -- node[m, above]{offer (SDP) over WebSocket} (4.3,-1.0);
  % 2 push
  \node[n] at (4.55,-1.4) {2};
  \draw[msg, draw=sheetRed] (4.3,-1.4) -- node[m, above]{\ct{apns-push-type: voip} · topic \ct{<bundle>.voip} · expiration 0} (10.1,-1.4);
  % 3 report
  \node[n] at (9.85,-1.8) {3};
  \node[act, draw=sheetRed, fill=sheetRed!7, anchor=west] at (10.25,-1.85) {\ct{reportNewIncomingCall} \textbf{now} $\to$ system call UI (even locked)};
  % 4 connect
  \node[n] at (12.75,-2.3) {4};
  \draw[msg] (13.0,-2.3) -- node[m, above, pos=0.3]{connect in parallel, fetch the offer} (4.3,-2.3);
  % 5 answer
  \node[n] at (10.35,-2.7) {5};
  \draw[msg, draw=sheetOrange] (10.1,-2.7) -- node[m, above]{\ct{CXAnswerCallAction}} (13.0,-2.7);
  \draw[msg] (13.0,-3.05) -- node[m, above, pos=0.25]{answer (SDP): \ct{setRemote(offer)} $\to$ \ct{createAnswer} $\to$ \ct{setLocal}} (0.9,-3.05);
  % 6 ICE
  \node[n] at (0.65,-3.5) {6};
  \draw[ice] (0.9,-3.5) -- node[m, above, pos=0.2]{ICE candidates, trickled both ways via signalling} (13.0,-3.5);
  \draw[<->, thick, draw=sheetBrown] (13.0,-3.85) -- node[m, above]{binding $\to$ public addr} (15.75,-3.85);
  \node[m, anchor=west, text=sheetBrown] at (13.1,-4.15) {TURN: relayed addr};
  % 7 media
  \node[n] at (0.65,-4.4) {7};
  \draw[med] (0.9,-4.4) -- node[m, above, pos=0.5]{connectivity checks $\to$ DTLS handshake $\to$ \textbf{SRTP media} (peer-to-peer, or via TURN)} (13.0,-4.4);
  % 8 audio
  \node[n] at (9.85,-4.85) {8};
  \draw[msg, draw=sheetOrange] (10.1,-4.85) -- node[m, above]{\ct{didActivate audioSession}} (13.0,-4.85);
  \node[m, anchor=west, align=left] at (13.1,-4.85) {start audio;\\fulfil the answer};
\end{tikzpicture}

\begin{multicols}{2}
\footnotesize\setstretch{1.0}\raggedright

\hd{WebRTC}
\begin{itemize}
  \item \textbf{Signalling is not specified} by WebRTC (JSEP, RFC 8829/9429): you move SDP and
        candidates over your own channel. \textbf{SDP} describes codecs, media sections, ICE ufrag/pwd
        and the \textbf{DTLS certificate fingerprint}. Offer/answer (RFC 3264)\unverified:
        \ct{createOffer} $\to$ \ct{setLocalDescription} $\to$ send; the peer
        \ct{setRemoteDescription} $\to$ \ct{createAnswer} $\to$ \ct{setLocalDescription} $\to$ send back.
  \item \textbf{ICE} (RFC 8445) gathers \textbf{candidates}: \emph{host} (LAN IP), \emph{server-reflexive}
        (public IP:port a \textbf{STUN} server saw, RFC 8489), \emph{relay} (a \textbf{TURN} server
        forwards, RFC 8656). Pairs are tested with STUN connectivity checks; the best working pair
        wins. \textbf{Trickle ICE} (RFC 8838)\unverified{} sends candidates as they are found.
  \item \textbf{NAT}: STUN + simultaneous sends punch holes in most NATs; a \emph{symmetric} NAT or a
        UDP-blocking firewall defeats it $\to$ TURN (UDP, TCP or TLS on 443). TURN costs you
        bandwidth — budget for it. Network change (Wi-Fi $\to$ cellular) $\to$ \textbf{ICE restart}.
  \item \textbf{Security}: media is always encrypted. \textbf{DTLS} handshake on the chosen path,
        checked against the SDP fingerprint; its keys feed \textbf{SRTP} (RFC 5764)\unverified. So a
        tampered signalling channel = MITM: protect it (TLS + auth). Data channels = SCTP over DTLS.
  \item \textbf{Codecs}: Opus audio; VP8 and H.264 mandatory for video. \textbf{Group calls}: mesh
        (N$-$1 uplinks each: small groups only) · \textbf{SFU} forwards streams (simulcast layers) · MCU mixes
        (server CPU).
  \item \textbf{On iOS}: no Apple WebRTC framework for native apps — ship libwebrtc
        (\ct{RTCPeerConnectionFactory}, \ct{RTCPeerConnection}, \ct{RTCIceServer}) or a vendor SDK.
        With CallKit set \ct{RTCAudioSession.useManualAudio = true} and enable audio only in
        \ct{didActivate}.
\end{itemize}

\hd{Example — PushKit $\to$ CallKit (the must-report path)}
\begin{lstlisting}[language=SwiftSheet]
let registry = PKPushRegistry(queue: .main)
registry.delegate = self; registry.desiredPushTypes = [.voIP]
func pushRegistry(_ r: PKPushRegistry,
    didUpdate c: PKPushCredentials, for t: PKPushType) {
  api.uploadVoIPToken(c.token) }          // not the alert token
func pushRegistry(_ r: PKPushRegistry,
    didReceiveIncomingPushWith p: PKPushPayload, for t: PKPushType,
    completion: @escaping () -> Void) {
  let from = p.dictionaryPayload["from"] as? String ?? "Unknown"
  let u = CXCallUpdate()
  u.remoteHandle = CXHandle(type: .generic, value: from)
  let id = calls.uuid(for: p.dictionaryPayload)  // server callId to UUID
  provider.reportNewIncomingCall(with: id, update: u) { _ in completion() }
  signalling.connect()                   // in parallel, not before
}
func provider(_ p: CXProvider, perform a: CXAnswerCallAction) {
  Task { await call.answer(); a.fulfill() } }   // fulfil once connected
\end{lstlisting}

\hd{Background}
VoIP push launches or wakes the app (even after a force-quit). During the call the \ct{audio} +
\ct{voip} background modes keep it running while audio flows; after hang-up it suspends as usual. The
old always-connected VoIP socket (\ct{setKeepAliveTimeout}) is deprecated — PushKit replaced it.

\hd{Over WebSocket, not WebRTC}
\textbf{SignalR} (.NET: hubs, falls back WebSocket $\to$ SSE $\to$ long polling) and \textbf{Socket.IO}
(Engine.IO: rooms, acks, starts with HTTP long polling then upgrades) are \emph{protocols on top} — a
plain WebSocket client can't talk to them; use their client library.

\columnbreak

\hd{CallKit + PushKit — the rules}
\begin{itemize}
  \item \textbf{\ct{CXProvider}} (\ct{CXProviderConfiguration}: \ct{supportsVideo},
        \ct{supportedHandleTypes}, \ct{maximumCallsPerCallGroup}, ringtone, icon) = \emph{system}
        $\to$ app: it reports calls (\ct{reportNewIncomingCall}, \ct{reportOutgoingCall(with:connectedAt:)},
        \ct{reportCall(with:endedAt:reason:)}) and delivers \emph{actions} to its delegate.
        \textbf{\ct{CXCallController}} = app $\to$ system: \ct{request(CXTransaction(action:))} with
        \ct{CXStartCallAction}, \ct{CXEndCallAction}, \ct{CXSetMutedCallAction},
        \ct{CXSetHeldCallAction}. Every action: do the work, then \ct{fulfill()} or \ct{fail()}.
        \ct{providerDidReset} $\to$ end everything.
  \item \textbf{The iOS 13 rule} (Apple, PushKit docs): linking the iOS 13+ SDK, every
        \ct{.voIP} push \textbf{must} be reported via \ct{reportNewIncomingCall}; \emph{``if you fail to
        report a call to CallKit, the system will terminate your app''}, and repeated failures may
        stop VoIP push delivery. So: report first, network second.
  \item No CallKit (e.g. a chat ``ping'') $\to$ no PushKit: use ordinary pushes + an NSE. Caller hung
        up before you connected $\to$ still report, then \ct{reportCall(\ldots reason: .remoteEnded)};
        can't reach your server $\to$ \ct{.failed}. Don't send more VoIP pushes to cancel — use your
        connection.
  \item \textbf{Push}: \ct{apns-push-type: voip}, \ct{apns-topic} = bundle ID + \ct{.voip},
        \ct{apns-expiration: 0} (a ring that arrives late is worse than none); token from
        \ct{didUpdate pushCredentials} — separate from the alert token.
  \item \textbf{Outgoing}: \ct{CXCallController.request} a \ct{CXStartCallAction} $\to$ your delegate's
        \ct{perform} starts signalling, \ct{fulfill()} $\to$
        \ct{reportOutgoingCall(with:startedConnectingAt:)} $\to$ \ldots\ct{connectedAt:} when media flows.
        Recents / Siri hand the app an \ct{INStartCallIntent}; it then requests the same action.
\end{itemize}

\hd{Audio session during a call}
Set \ct{.playAndRecord} + mode \ct{.voiceChat} (echo cancellation; \ct{.videoChat} for video) early,
but \textbf{don't activate it yourself}: CallKit activates it (with call priority) and calls
\ct{provider(\_:didActivate:)} — start audio I/O there, stop in \ct{didDeactivate}. Hold / a second
call arrives as \ct{CXSetHeldCallAction}; mute as \ct{CXSetMutedCallAction} $\to$ disable the audio track.

\hd{Traps}
\begin{itemize}
  \trap{Doing network work before \ct{reportNewIncomingCall} — the app is killed on iOS 13+.}
  \trap{Using VoIP pushes for messages ``to wake the app'' — not allowed; NSE instead.}
  \trap{\ct{setActive(true)} yourself / starting audio before \ct{didActivate} — silent call.}
  \trap{``STUN is enough'' — some users need TURN; no TURN = calls that never connect.}
  \trap{Thinking WebRTC defines signalling, or that SRTP keys travel in SDP (only the fingerprint does).}
\end{itemize}

\hd{Remember}
\textbf{Push $\to$ report $\to$ connect $\to$ answer $\to$ ICE $\to$ DTLS $\to$ SRTP $\to$ didActivate.}


\end{multicols}

\noindent{\footnotesize\color{sheetGrey}\textit{Related:} websockets-realtime (signalling transport,
reconnect) · push-notifications (APNs, NSE) · background-execution (modes) · avfoundation-camera-audio
(categories, interruptions) · crypto-cryptokit (DTLS keys) · system-design-sdk-messenger}

\end{document}
